% !TEX root = resume.tex
% !TeX program = lualatex
% !BIB program = biber
% !TEX bibfile = file.bib

%----------------------------------------------------------------------------------------
%	PACKAGES AND OTHER DOCUMENT CONFIGURATIONS
%----------------------------------------------------------------------------------------
% 
\documentclass[
	10pt, % Default font size, can be between 8pt and 12pt
]{RESUME}

%-----------------------
% Make PDF Dark

% \pagecolor[rgb]{0.5,0.5,0.5} % Gray backgrond

% \color[rgb]{0.5,0.5,0.5} %grey
%-----------------------

\columnratio{0.55, 0.45} % Widths of the two columns, specified here as a ratio summing to 1 to correspond to percentages; adjust as needed for your content 

% Headers and footers can be added with the following commands: \lhead{}, \rhead{}, \lfoot{} and \rfoot{}
% Example right footer:
%\rfoot{\textcolor{headings}{\sffamily Last update: \today. Typeset with Xe\LaTeX}}

%----------------------------------------------------------------------------------------

\begin{document}

\begin{paracol}{2} % Begin two-column mode

	%----------------------------------------------------------------------------------------
	%	YOUR NAME AND CURRICULUM VITAE TITLE
	%----------------------------------------------------------------------------------------

	\parbox[][0.11\textheight][c]{\linewidth}{ % Box to hold your name and CV title; change the fixed height as needed to match the colored box to the right
		\centering % Horizontally center text

		{\sffamily\Huge Atanu Kumar Biswas} % Your name

		\medskip % Vertical whitespace

		{\cursivefont\Huge\textcolor{headings}{Cyber Security Intern}}

		\vfill % Push content to the top of the box

	}


	%----------------------------------------------------------------------------------------
	%	EDUCATION
	%----------------------------------------------------------------------------------------

	\section{Education}

	% Each qualification entry is added with a \qualificationentry command. Below is an empty one to use as a template:

	%\qualificationentry
	%	{} % Duration
	%	{} % Degree
	%	{} % Honors, achievements or distinctions (e.g. first class honors)
	%	{} % Department
	%	{} % Institution

	% All 5 parameters must be supplied but any can be empty if you don't need them

	%------------------------------------------------

	\begin{supertabular}{r l} % Start a table with two columns, the table will ensure everything is aligned

		%------------------------------------------------

		\qualificationentry
		{2016 - 2022} % Duration
		{Bachelor of Science} % Degree
		{} % Honors, achievements or distinctions (e.g. first class honors)
		{Computer Science and Engineering } % Department
		{Chittagong University of Engineering \& Technology} % Institution
		%------------------------------------------------

	\end{supertabular}




	%----------------------------------------------------------------------------------------
	%	MAJOR RESEARCH PROJECT
	%----------------------------------------------------------------------------------------

	\section{Undergraduate Research}

	 {\raggedright\textbf{"A distributed and transparent food supply chain management system using blockchain and IoT technologies"}\par}

	\medskip % Vertical whitespace
	In this project I worked on food supply chain System. As we know that food supply chain
	is the most crucial supply chain on our planet. Existing food supply
	chain systems do not investigate features like food origin traceability, timestamps of food
	processing, a proper money distribution facility among the parties, and a customer feedback
	evaluation system at the same time. To make the food supply we've used IoT and
	blockchain technologies. Three different blockchains
	are developed for customers, warehouses, and farmers. Customer blockchain is public and
	anyone can participate in this. Other blockchains are private and require registration for
	participation. Our experimental results based on sensors and blockchain data show that our
	proposed scheme improves 96.83\% of the existing scheme's time.
	\medskip % Extra vertical whitespace before the next section

	%----------------------------------------------------------------------------------------
	%	WORK EXPERIENCE
	%----------------------------------------------------------------------------------------

	\section{Work Experience}

	% Each job is added with a \jobentry command. Below is an empty one to use as a template:

	%\jobentry
	%	{} % Duration
	%	{} % FT/PT (full time or part time)
	%	{} % Employer
	%	{} % Job title
	%	{} % Description

	% All 5 parameters must be supplied but any can be empty if you don't need them

	%------------------------------------------------

	\jobentry
	{Aug, 2019 - Sep, 2019} % Duration
	{FT} % FT/PT (full time or part time)
	{Contessa Solutions \& Consultants Ltd.} % Employer
	{ Student Intern } % Job title
	{As part of this internship, we learned how data center works. We created some virtual machine using
		\textit{VMware ESXi} formerly known as ESX is a bare metal hypervisor. We learn little bit of network
		virtualization. We created some virtual storage to test how data center works.} % Description

	%------------------------------------------------
	%------------------------------------------------

	\jobentry
	{2016 - 2017} % Duration
	{FT} % FT/PT (full time or part time)
	{The NASA Human Exploration Rover Challenge (HERC)} % Employer
	{Team Leader (Software Team)} % Job title
	{As part of this project, we learned \textit{Arduino, RaspbarryPi, GPS, Motor Operation, Robotic Arm
			Operation, Different type or sensor ( pH, moisture, sonar, IR }. We also to build \textit{image processing}
		tool that
		can detect colored ( for example red) ball and follow the path accordingly. We communicate with the robot
		using \textit{radio} communication and \textit{Bluetooth}. We used \textit{Socket Programming} to
		communicate with server. Our team selected for the competition.} % Description

	%------------------------------------------------
	%------------------------------------------------

	\jobentry
	{Jun, 2016 - Aug, 2016} % Duration
	{FT} % FT/PT (full time or part time)
	{Robotic Competition} % Employer
	{Team Leader (Robot Racing Team CSE, CUET)} % Job title
	{As part of this promotion, we learned hot to build a robot that can detect black line on white surface
		and follow the path. We programmed the robot using Arduino and IR sensor. Our task is to follow the path
		and stop at the destination as fast as possible.} % Description

	%------------------------------------------------


	%------------------------------------------------
	\break
	%------------------------------------------------

	%----------------------------------------------------------------------------------------
	%	REFERENCES
	%----------------------------------------------------------------------------------------
	\section{References}

	%\textit{References available on request} % Uncomment if you'd rather not include references and remove the section below

	%------------------------------------------------

	% This section is laid out using a table. A \tableentry command adds lines with the following parameters:

	%\tableentry{Heading}{Content}{spaceafter}
	% All 3 parameters must be supplied but any can be empty if you don't need them
	% A "spaceafter" value in the third parameter will add some vertical space -- this is to be used between headings, leave it empty for no extra space

	%------------------------------------------------

	\begin{supertabular}{r l} % Start a table with two columns, the table will ensure everything is aligned

		%------------------------------------------------

		\tableentry{}{\textbf{Pranto Biswas}}{}
		\tableentry{}{\href{https://www.linkedin.com/in/pranxol/}{\textit{Software Engineer}}}{spaceafter}
		\tableentry{Position}{Software Engineer L2}

		\tableentry{Organization}{\href{https://www.linkedin.com/company/681429/}{\textit{Enosis Solutions}}}{spaceafter}
		\tableentry{Email}{\href{mailto:pranto.biswas@enosisbd.com}{pranto.biswas@enosisbd.com} (Work)}{}
		\tableentry{Phone}{+88 01521-221409 (Work)}{}

		%------------------------------------------------

		\\ % Additional vertical whitespace between the references


	\end{supertabular}

	\medskip % Extra vertical whitespace before the next section

	%----------------------------------------------------------------------------------------

	\switchcolumn % Switch to the second (right) column

	%----------------------------------------------------------------------------------------
	%	COLORED CONTACT DETAILS BOX
	%----------------------------------------------------------------------------------------

	\parbox[top][0.11\textheight][c]{\linewidth}{ % Box to hold the colored box; change the fixed height as needed to match the box to the left
		\colorbox{shade}{ % Create colored box and specify background color
			\begin{supertabular}{@{\hspace{3pt}} p{0.05\linewidth} | p{0.775\linewidth}} % Start a table with two columns, the table will ensure everything is aligned
				\raisebox{-1pt}{\faHome} & Padmapukur, Koyra, Khulna, 9290 \\ % Address

				\raisebox{-1pt}{\faPhone} & +88 01907-956568 \\ % Phone number

				\raisebox{-1pt}{\small\faEnvelope} & \href{mailto:BKumarAtanu@zohomail.com}{bkumaratanu@zohomail.com} \\ % Email address
				\raisebox{-1pt}{\small\faLinkedin} & \href{https://www.linkedin.com/in/bkumaratanu/}{linkedin.com/in/BKumarAtanu}\\ % Linkedin

				\raisebox{-1pt}{\small\faGlobe} & \href{https://bkumaratanu.github.io/BKumarAtanu/}{bkumaratanu.github.io/BKumarAtanu}\\ % Website

				\raisebox{-1pt}{\small\faGithub} & \href{https://github.com/BKumarAtanu}{github.com/BKumarAtanu}\\ % Github

				% \raisebox{-1pt}{\small\faGitlab} & \href{https://gitlab.com/BKumarAtanu}{gitlab.com/BKumarAtanu}\\ % Gitlab

				%\raisebox{-1pt}{\faGithub} & \href{https://github.com/username}{https://github.com/username} \\ % GitHub profile
				%\raisebox{-1pt}{\faLinkedinSquare} & \href{https://www.linkedin.com/in/username}{https://www.linkedin.com/in/username} \\ % LinkedIn profile
				% See fontawesome.pdf in the Fonts folder for all icons you can use
			\end{supertabular}
		}
		\vfill % Push content to the top of the box
	}



	%----------------------------------------------------------------------------------------
	%	AWARDS
	%----------------------------------------------------------------------------------------

	\section{Competitions}

	% This section is laid out using a table. A \tableentry command adds lines with the following parameters:

	%\tableentry{Heading}{Content}{spaceafter}
	% All 3 parameters must be supplied but any can be empty if you don't need them
	% A "spaceafter" value in the third parameter will add some vertical space -- this is to be used between headings, leave it empty for no extra space

	%------------------------------------------------

	\begin{supertabular}{r l} % Start a table with two columns, the table will ensure everything is aligned

		%------------------------------------------------

		\tableentry{2016}{\textbf{Roborace}}{}
		\tableentry{}{\textit{\href{https://www.cuet.ac.bd/}{CUET}}}{spaceafter}

		%------------------------------------------------

		\tableentry{2017}{\textbf{NASA Mars Exploration Rovers}}{}
		\tableentry{}{\textit{\href{https://www.nasa.gov/learning-resources/nasa-human-exploration-rover-challenge/}{NASA Human Expl. Rover Challenge (HERC)}}}{spaceafter}

		%------------------------------------------------

	\end{supertabular}

	%----------------------------------------------------------------------------------------
	%	SKILLS DESCRIPTION
	%----------------------------------------------------------------------------------------

	\section{Programming Skills}

	\begin{supertabular}{r l} % Start a table with two columns, the table will ensure everything is aligned

		%------------------------------------------------

		\tableentry{Beginner}{R, Go, API,}{}
		\tableentry{}{HTML, CSS, JavaScript, GNU Octave }{spaceafter}

		%------------------------------------------------
		\tableentry{Intermediate}{C, C++, SQL, Git}{}
		\tableentry{}{Shell Scripting}{spaceafter}

		%------------------------------------------------

		\tableentry{Expert}{Python}{spaceafter}

		%------------------------------------------------
	\end{supertabular}


	%----------------------------------------------------------------------------------------
	%	COMPUTER SKILLS
	%----------------------------------------------------------------------------------------

	\section{Technical Proficiency}

	% This section is laid out using a table. A \tableentry command adds lines with the following parameters:

	%\tableentry{Heading}{Content}{spaceafter}
	% All 3 parameters must be supplied but any can be empty if you don't need them
	% A "spaceafter" value in the third parameter will add some vertical space -- this is to be used between headings, leave it empty for no extra space

	%------------------------------------------------

	\begin{supertabular}{r l} % Start a table with two columns, the table will ensure everything is aligned

		%------------------------------------------------

		\tableentry{Beginner}{Inkscape, Libreoffice }{spaceafter}

		%------------------------------------------------
		\tableentry{Intermediate}{Anaconda, Calibre}{}
		\tableentry{}{\LaTeX, Markdown}{}
		\tableentry{}{Microsoft Windows}{}
		\tableentry{}{Computer Hardware \& Support}{spaceafter}

		%------------------------------------------------

		\tableentry{Expert}{Linux OS}{spaceafter}

		%------------------------------------------------

	\end{supertabular}


	%----------------------------------------------------------------------------------------
	%   CERTFICATIONS	
	%----------------------------------------------------------------------------------------


	\section{Online Certifications}

	%------------------------------------------------
	\subsection{\textbf{Machine Learning}}
	\begin{itemize}

		\item[\faCertificate] \href{https://www.datacamp.com/statement-of-accomplishment/course/89efb9c64a2ba11dbe758718478b1e6584482f53?raw=1}
		      {\textcolor{black}{\textit{Supervised Learning with scikit-learn}}}

		\item[\faCertificate] \href{https://www.datacamp.com/statement-of-accomplishment/course/5a11a956cbbf72f58df34ebb4bcb8953d50827da?raw=1}
		      {\textcolor{black}{\textit{Unsupervised Learning in Python}}}

		\item[\faCertificate] \href{https://www.datacamp.com/statement-of-accomplishment/course/4ecb52a63fbb1b3a1ca42c34374826979cd1cb48?raw=1}
		      {\textcolor{black}{\textit{Machine Learning with Tree-Based Models in Python}}}

		\item[\faCertificate] \href{https://www.datacamp.com/statement-of-accomplishment/course/66403a5d2f4b1113615c4332ff47006d7d6b273d?raw=1}
		      {\textcolor{black}{\textit{Hypothesis Testing in Python}}}

		\item[\faCertificate] \href{https://mega.nz/file/olQSCLgR\#MNku1sr157kYUQNx9KW2Ajq9rqzuz5S4D2XCS38LaSc}
		      {\textcolor{black}{\textit{Machine Learning and Artificial Intelligence Fundamentals with Python}}}

	\end{itemize}

	\subsection{\textbf{Data Manupulation}}
	\begin{itemize}

		\item[\faCertificate] \href{https://www.datacamp.com/statement-of-accomplishment/course/fbc94a874eb1ade1ed86b66c19fd32b3c5077f2a?raw=1}
		      {\textcolor{black}{\textit{Sampling in Python}}}

		\item[\faCertificate] \href{https://www.datacamp.com/statement-of-accomplishment/course/62a92ec8ef730881263a99b128064cc771ca7519?raw=1}
		      {\textcolor{black}{\textit{Data Manipulation with pandas}}}

		\item[\faCertificate] \href{https://www.datacamp.com/statement-of-accomplishment/course/55a0802fb62bd726dbdb6906505f54d98f8f1c6c?raw=1}
		      {\textcolor{black}{\textit{Manipulating Time Series Data in Python}}}
	\end{itemize}
	% Customized bullets
	\subsection{\textbf{Data Visualization}}

	\begin{itemize}

		\item[\faCertificate] \href{https://www.datacamp.com/statement-of-accomplishment/course/17c9ab4dc4eb990e81e2b675f3d6375691e66614?raw=1}
		      {\textcolor{black}{\textit{Data Visualization with Matplotlib}}}

		\item[\faCertificate] \href{https://www.datacamp.com/statement-of-accomplishment/course/0193c3a65db0fe46cd4d1dda686015f6964744bf?raw=1}
		      {\textcolor{black}{\textit{Data Visualization with Seaborn}}}
	\end{itemize}


	% \item[\faCertificate] \href{}
	%       {\textcolor{black}{\textit{}}}
	\subsection{\textbf{Scripting}}

	\begin{itemize}

		\item[\faCertificate] \href{https://www.datacamp.com/statement-of-accomplishment/course/af330fd7f2c7b67b78329c57e31171f67dd2daba?raw=1}
		      {\textcolor{black}{\textit{Introduction to Bash Scripting}}}

	\end{itemize}

	\break
	%----------------------------------------------------------------------------------------
	%	COMMUNICATION SKILLS
	%----------------------------------------------------------------------------------------

	\section{Communication Skills}

	% This section is laid out using a table. A \tableentry command adds lines with the following parameters:

	%\tableentry{Heading}{Content}{spaceafter}
	% All 3 parameters must be supplied but any can be empty if you don't need them
	% A "spaceafter" value in the third parameter will add some vertical space -- this is to be used between headings, leave it empty for no extra space

	%------------------------------------------------

	\begin{supertabular}{r l} % Start a table with two columns, the table will ensure everything is aligned

		%------------------------------------------------

		\tableentry{Languages}{Bangla, English}{}
		\tableentry{Bangla}{\faStar\faStar\faStar\faStar \faStarO }{}
		\tableentry{English}{\faStar \faStar \faStar \faStarO \faStarO}{}
		% \tableentry{Hindi}{\faStar \faStar \faStarHalfFull \faStarO \faStarO}{spaceafter}

		%------------------------------------------------

		\tableentry{Team Work}{\faStar \faStar \faStar \faStarHalfFull \faStarO} {}
		\tableentry{}{I've worked with \textbf{4} teams already}{spaceafter}

		%------------------------------------------------

		\tableentry{Team Leading}{\faStar \faStar \faStar \faStarHalfFull \faStarO} {}
		\tableentry{}{I've leaded \textbf{2} teams already}{spaceafter}

	\end{supertabular}

	%------------------------------------------------
	%------------------------------------------------

	%----------------------------------------------------------------------------------------
	%	PUBLICATIONS
	%----------------------------------------------------------------------------------------
	\section{Publications}

	%------------------------------------------------

	\textbf{Atanu Kumar Biswas} (2022) \\
	\medskip % Vertical whitespace
	\textbf{Mahfuzulhoq Chowdhury} (2022).\\
	A distributed and transparent food supply chain management system using blockchain and IoT technologies. \textit{International Journal of Society Systems Science, Vol. 14, No. 1}.

	\medskip % Vertical whitespace

	%------------------------------------------------

	% As an alternative to a long-form publication list, you can create a shorter summary using only DOI values and years.

	% Example \doipublication{} command to add another publication:

	%\doipublication{Year}{DOI}{firstauthor}{spaceafter}

	% All four parameters are required (can be empty though)
	% A value of "firstauthor" in the third parameter will output the DOI in bold
	% A "spaceafter" value in the fourth parameter will add some vertical space -- this is to be used between years

	%------------------------------------------------

	\subsection{Publications by DOI}

	\begin{supertabular}{r l} % Start a table with two columns, the table will ensure everything is aligned

		%------------------------------------------------

		\doipublication{2022}{10.1504/IJSSS.2022.128099}{firstauthor}{spaceafter}

		%------------------------------------------------

		& \textit{First author publications in} \textbf{first place}\\

		%------------------------------------------------

	\end{supertabular}

	\medskip % Extra whitespace before the next section

	%----------------------------------------------------------------------------------------

\end{paracol} % End two-column mode

%----------------------------------------------------------------------------------------

\end{document}
